\section{More Details On the Setup}

\subsection{LLM-Based Evaluator.} 
\label{cohen's kappa}
Cohen’s Kappa results are detailed in Table \ref{tab:cohen's kappa}, illustrating the level of concordance between expert annotators and LLM-evaluators. Our result reveals a Cohen's Kappa value of 0.520 for the 0-0.5-1 scale and 0.444 for the 1-2-3-4-5 scale. While the Cohen's Kappa value of 0.520 only indicates a ``weak agreement'' with humans \cite{cohen's_kappa}, the 0-0.5-1 scale demonstrates stronger agreement compared to the 1-2-3-4-5 scale. Additionally, the evaluation prompts used for both the 0-0.5-1 and 1-2-3-4-5 scales are provided in Table \ref{tab:gpt_based_evaluator_app}. 

\begin{table}[!t]
    \centering
    \small
    \begin{tabular}{ccc}
    \toprule
     & \textbf{0-0.5-1} & \textbf{1-2-3-4-5} \\
    \midrule
    Expert Annotators & 0.520 & 0.444 \\
    \bottomrule
    \end{tabular}
    \caption{Cohen's Kappa between two LLM-evaluators \wrt expert annotations.}
    \label{tab:cohen's kappa}
\end{table}






\subsection{Prompt for Evaluating Visual Context Question}
Prompts used to rank reference papers across multiple LMMs are illustrated in Table \ref{tab:locality prompt}.
\begin{table*}[!t]
\centering
\small
\renewcommand\tabcolsep{1.0pt} % column space
\begin{tabular}{cp{11.5cm}}
    \toprule
    \textbf{Model} & \textbf{Prompt} \\
    \midrule
    \multirow{3}{*}{\parbox{4cm}{\centering Yi-VL-34B\\DeepSeek-VL}}
    &  Answer the question from the figure and the reference papers provided only: \{question\}
    
    Additionally, rerank the following reference papers according to their relevance to this question. Each reference paper consists of an S2\_id, a title, and an abstract.
    
    \{paper\_cluster\}
    
    Format your answer as a python dictionary with keys "question", "answer", and "rank". "rank" should be a list of S2\_id. 
    
    If no relevant reference papers are provided, return an empty list for "rank". Note: The "rank" list should only include **question-relevant** reference papers. Do not include irrelevant ones.\\
    \midrule
    \multirow{3}{*}{\parbox{4cm}{\centering InternVL-Chat-1.1V}}
    &  You are given a figure, a question, and some paper candidates of titles and abstracts. Your task is to answer the question based on the figure information, then order the paper candidates that I provide to you so that the paper that is more relevant to the question comes first in the list.
    Provide your answer at the end in a json file of this format using S2\_id only:\{"ranking":  [""] \}. Make sure the responded list is in a valid format and that it only contains the S2\_id. Do not include the title or abstract in the answer list. 

    <question>
    \{question\}
    </question>

    <paper candidates>
    \{paper\_cluster\}
    </paper candidates>\\
    \midrule
    \multirow{3}{*}{\parbox{4cm}{\centering LLaVa-1.6\\Qwen-VL}}
    & Answer the question from the figure and the reference papers provided only: \{question\}
    
    Additionally, rerank the following reference papers according to their relevance to this question. Each reference paper consists of an S2\_id, a title, and an abstract.
    
    \{paper\_cluster\}
    
    Format your answer as a python dictionary with keys "question", "answer", and "rank". "rank" should be a list of S2\_id. 
    
    If no relevant reference papers are provided, return an empty list for "rank".
    \\
    \midrule
    \multirow{8}{*}{\parbox{4cm}{\centering GPT-4o\\GPT-4V(ision)\\Gemini-Pro-Vision-1.0\\Claude-3-Haiku\\Claude-3-Sonnet\\Claude-3-Opus}}
    & You are given a figure, a question, and a list of paper candidates of titles and abstracts. Your task is to answer the question based on the figure information and then re-rank the list of paper candidates I provided to you. 

        Provide your answer at the end in a json format using the S2\_id only: \{"ranking":  []\}.
        Only include papers that are relevant. Do not include papers that are irrelevant. Make sure the answer list is properly formatted.

        <question>
        \{question\}
        </question>

        <paper candidates>
        \{paper\_cluster\}
        </paper candidates>
    \\
    \bottomrule
\end{tabular}
\caption{Prompts used to rank reference papers across multiple LMMs.}
\label{tab:locality prompt}
\end{table*}


\subsection{Prompt for Answering Reference-based Question}
\label{appx:prompt_for_combined_questions}
Prompts used to answer reference-based questions are illustrated in Table \ref{tab:detail prompt}.

\begin{table*}[!t]
\centering
\small
\renewcommand\tabcolsep{1.0pt} % column space
\begin{tabular}{cp{11.5cm}}
    \toprule
    \textbf{Stage} & \textbf{Prompt} \\
    \midrule
    \multirow{3}{*}{\parbox{4cm}{\centering Answers from text chunk}}
    &  Answer the below question about a scientific paper. The question is composed of 2 parts, and the second part of the question can be answered from the paper. I will provide you with only a chunk of a paper. 
        
Explain your reasoning. Append the answer at the end of the response in a json format \{``answer'': ``''\}. You should answer the question in a short-answer form. Do not provide long answers. If you do not know the answer, respond with \{``answer'': ``I don't know''\}
        
<QUESTION>
\{question\}
</QUESTION>

<CHUNK>
\{chunk\}
</CHUNK>\\
    \midrule
    \multirow{3}{*}{\parbox{4cm}{\centering Answer aggregation}}
    &  I will provide you with a set of answer candidates for a question. Aggregate the information from all the candidates and give me one single answer. Note that if one answer candidate is `I don't know', you can ignore it.
    Answer the question based on the answer candidates and summarize the final answer into a short answer.
    <QUESTION>
    \{question\}
    </QUESTION>
    \{answer\_candidate\_list\}\\

    \bottomrule
\end{tabular}
\caption{Prompts used to generate and aggregate answers from a text chunk.}
\label{tab:detail prompt}
\end{table*}

\subsection{Prompt for Answer Evaluation}
Prompts used to retrieve answers from each text chunk and aggregate answers are illustrated in Table \ref{tab:gpt_based_evaluator_app}.
\begin{table*}[!t]
\centering
\small
\renewcommand\tabcolsep{1.0pt} % column space
\begin{tabular}{cp{11.5cm}}
    \toprule
    \textbf{Evaluator} & \textbf{Prompt} \\
    \midrule
    \multirow{8}{*}{\parbox{4cm}{\centering LLM-based Evaluator \\(0-0.5-1 setting)}}
    &  I am testing a model's performance on open-ended questions. I want you to help me in checking to see if the candidate answer has the same meaning as the reference answer for a given question. If you think the reference answer and the candidate answer have the same meaning, respond \{``selection'': ``1''\}; otherwise, respond by \{``selection'': ``0''\}. If you think the candidate is partially correct, respond by \{``selection'': ``0.5''\}. If the answer is ``I don't know,'' rate it to 0.

    <QUESTION>
    \{question\}
    </QUESTION>

    <REFERENCE>
    \{reference\}
    </REFERENCE>

    <CANDIDATE>
    \{candidate\}
    </CANDIDATE>
\\
\midrule

\multirow{28}{*}{\parbox{4cm}{\centering LLM-based Evaluator \\(1-2-3-4-5 setting)}}
    &  I am testing a model's performance on open-ended questions. I want you to help me in checking to see if the candidate answer has the same meaning as the reference answer for a given question. 
    
    Rate the candidate answer from 1, 2, 3, 4, and 5, where 1 means the candidate is the least similar to the reference answer and 5 means the candidate matches to the reference answer perfectly. 
    Respond by \{``selection'': ``''\}. If the candidate answer is ``I don't know,'' rate it to 1.

    Here's some examples you can consider:
    
    Question: Why transformer is better than RNN?
    
    Reference: Parallel computation
    
    Candidate: Computation
    
    Rating: 3

    Question: What's the major advantage of using ALiBi positional embedding?
    
    Reference: Effectively handle sequences of varying lengths, particularly beneficial for very long sequences
    
    Candidate: It has more freedom to handle input
    
    Rating: 2

    Question: What's the model's performance on GSK8K dataset? 
    
    Refernece: 65.65\%
    
    Candidate: 44.56\%
    
    Rating: 1

    Question: What specific method does this paper propose to solve LLM searching problem?
    
    Reference: MCTS
    
    Candiate: Monte Carlo Tree Search is proposed in this paper to solve searching when using decomposed prompting method.
    
    Rating: 5

    Question: How does the performance change when we switch from CoT to ToT in prompting? 
    
    Reference: Accuracy from 23.50\% to 32.87\%
    
    Candidate: slightly increase
    
    Rating: 4

    <QUESTION>
    \{question\}
    </QUESTION>

    <REFERENCE>
    \{reference\}
    </REFERENCE>

    <CANDIDATE>
    \{candidate\}
    </CANDIDATE>
\\
    \bottomrule
\end{tabular}
\caption{Prompts used to evaluate answers generated by LLMs.}
\label{tab:gpt_based_evaluator_app}
\end{table*}

\subsection{Prompt for Reference-based Question Generation}
\label{prompt_detail}
We list our prompt for reference-based question generation in Table \ref{tab:E.5}.


\begin{table*}[!t]
\centering
\small
\renewcommand\tabcolsep{1.0pt} % column space
\begin{tabular}{cp{11.5cm}}
    \toprule
    \textbf{Model} & \textbf{Prompt} \\
    \midrule

    \multirow{18}{*}{\parbox{4cm}{\centering Reference-based Question Generation Prompt}}
    &  Generate 1 short answer question based on the paper's full content below. You should follow the reasoning type of \{reasoning\_type\}, with the definition \{reasoning\_description\}. The short answer question should be as hard as possible, and focus on a single detail from the paper. The target audience of the short answer question is an expert in the field of natural language processing. The question should be hard for GPT-4 to answer. The answer to the question should be short and must be answerable from the content of the paper.

Here are some requirements: 

<REQUIREMENTS>
[Question] should make sense and can be answered from the paper's full text.
[Answer] should be directly answering the question you generated.
[Explanation] should explain why the answer correctly answers the question.
[Evidence] should be from the original content from the paper content. This should be an excerpt from the input paper that supports your answer. 

</REQUIREMENTS>

Append the answer at the end of your response in a json-like format:

\{``question'': ``'', 

``answer'': ``'', 

``explanation'': ``'', 

``evidence'':``''\}

<PAPER FULL CONTENT>

{full\_text}

</PAPER FULL CONTENT>\\

    \bottomrule
\end{tabular}
\caption{Prompt for reference-based question generation.}
\label{tab:E.5}
\end{table*}





\subsection{Model Parameters for Answering Visual Context Question}
Model parameters for ranking reference papers from a paper cluster are shown in Table \ref{tab:model_param_locality}.
\begin{table*}[!t]
\centering
\small
% \renewcommand\tabcolsep{1.0pt} % column space
\begin{tabular}{cl}
    \toprule
    \textbf{Model} & \textbf{Generation Setup} \\
    \midrule
    LLaVa-1.6 & model = \texttt{llava-v1.6-mistral-7b}, temperature = 0.1, max\_tokens = 8192\\
    \midrule
    Yi-VL-6B & model = \texttt{Yi-VL-6B}, temperature = 0.1, max\_tokens = 8192\\
    \midrule
    DeepSeek-VL & model = \texttt{deepseek-vl-7b-chat}, temperature = 0.1, max\_tokens = 8192\\
    \midrule
    InternVL-Chat-1.1V & model = \texttt{InternVL-Chat-Chinese-V1-1}, temperature = 0.1, max\_tokens = 768 \\
    \midrule
    Qwen-VL-Plus & model = \texttt{qwen-vl-plus}, seed = 1234, max\_tokens = 6000\\
    \midrule
    GPT-4o 
    & model = \texttt{gpt-4o}, temperature = 0.1, max\_tokens = 4096\\
    \midrule
    GPT-4V(ision) 
    & model = \texttt{gpt-4-turbo}, temperature = 0.1, max\_tokens = 4096\\
    \midrule
    Gemini-Pro-Vision-1.0 
    & model = \texttt{gemini-pro-vision}, temperature = 0.1, max\_tokens = 4096\\
    \midrule
    Claude-3-Haiku 
    & model = \texttt{claude-3-haiku-20240307}, temperature = 0.1, max\_tokens = 4096\\
    \midrule
    Claude-3-Sonnet 
    & model = \texttt{claude-3-sonnet-20240229}, temperature = 0.1, max\_tokens = 4096\\
    \midrule
    Claude-3-Opus 
    & model = \texttt{claude-3-opus-20240229}, temperature = 0.1, max\_tokens = 4096\\
    \bottomrule
\end{tabular}
\caption{Parameters of various LMMs in evaluating visual context questions.}
\label{tab:model_param_locality}
\end{table*}







\subsection{Model Parameters for Answering Reference-based Question}
Model parameters for answering reference-based questions are exhibited in Table \ref{tab:model_param_detail}.
\begin{table*}[!t]
\centering
\small
\begin{tabular}{cl}
    \toprule
    \textbf{Model} & \textbf{Generation Setup} \\
    \midrule
    Llama-3-70B & temperature = 0.1, max\_token = 10,000  \\
    \midrule
    Mistral-7B & temperature = 0.1, max\_token = 40,000\\
    \midrule
    PaLM-2 & temperature = 0.1, max\_token = 40,000 \\
    \midrule
    Gemma & temperature = 0.1, max\_token = 12,000\\
    \midrule
    DBRX & temperature = 0.1, max\_token = 40,000 \\
    \midrule
    Command R+ & temperature = 0.1, max\_token = 200,000 \\
    \midrule
    GPT-4 
    & model = \texttt{gpt-4-0125-preview}, temperature = 0.1, max\_tokens = 200,000\\
    \midrule
    Gemini-Pro-1.0 
    & model = \texttt{gemini-1.0-pro}, temperature = 0.1, max\_tokens = 40,000\\
    \midrule
    Claude-3-Haiku 
    & model = \texttt{claude-3-haiku-20240307}, temperature = 0.1, max\_tokens = 250,000\\
    \midrule
    Claude-3-Sonnet 
    & model = \texttt{claude-3-sonnet-20240229}, temperature = 0.1, max\_tokens = 250,000\\
    \midrule
    Claude-3-Opus 
    & model = \texttt{claude-3-opus-20240229}, temperature = 0.1, max\_tokens = 250,000\\
    \bottomrule
\end{tabular}
\caption{Parameters of various LLMs in evaluating reference-based questions.}\label{tab:model_param_detail}
\end{table*}




